% 1.1 — L'organisme d'accueil
\section{L'organisme d'accueil : WePwn}
\label{sec:wepwn}

\begin{tcolorbox}[enhanced,colback=wepwnclair,colframe=wepwn,boxrule=0pt,leftrule=4pt,arc=0pt,
                  left=10pt,right=10pt,top=8pt,bottom=8pt]
\begin{minipage}[c]{0.2\linewidth}
  \centering\includegraphics[width=0.85\linewidth]{figures/logos/wepwn.png}
\end{minipage}\hfill
\begin{minipage}[c]{0.76\linewidth}
  {\sffamily\Large\bfseries\color{wepwnfonce} WePwn}\par\smallskip
  {\sffamily\small Sécurité offensive · programmes de \textit{bug bounty} · divulgation
  coordonnée de vulnérabilités · outils de sécurité}
\end{minipage}
\end{tcolorbox}

\subsection{Présentation}
WePwn est une entreprise marocaine de sécurité offensive. Son activité d'origine est le
\textit{bug bounty} : une entreprise ouvre un programme, des chercheurs en sécurité
indépendants testent ses services, et chaque vulnérabilité validée est récompensée. WePwn
se présente comme la mise en relation entre des entreprises et des chercheurs, afin que
les premières soient averties des vulnérabilités identifiées chez elles~\cite{wepwn-site}.
Selon l'organisme, elle a été la première plateforme de \textit{bug bounty} au
Maroc\footnote{Affirmation de l'organisme d'accueil, reprise telle quelle.}~; les
programmes qu'elle héberge rémunèrent les chercheurs en dirhams. Un portail dédié permet
aux chercheurs de rejoindre ces programmes~\cite{wepwn-access}.

Autour de cette activité, WePwn propose un service de divulgation de vulnérabilités géré
de bout en bout~\cite{wepwn-services} : tri des signalements, évaluation de la sévérité,
validation, rapport détaillé, conseils de correction et pilotage du programme dans la
durée. Ce métier --- recevoir des attaques réelles, les reproduire, en mesurer l'impact
et accompagner la correction --- a naturellement conduit l'entreprise vers un second axe :
les \emph{pratiques} et les \emph{outils} de sécurité. \demo{} appartient à cet axe.

\subsection{Fiche signalétique}

\begin{table}[htbp]
\centering
\caption{Fiche signalétique de WePwn}
\label{tab:fiche-wepwn}
\begin{tabularx}{0.92\linewidth}{@{}L{4.2cm}Y@{}}
\toprule
Raison sociale & WePwn \\
Domaine & Cybersécurité --- sécurité offensive \\
Activités & Programmes de \textit{bug bounty} ; divulgation de vulnérabilités gérée ;
            pratiques et outils de sécurité \\
Fondateur & M. Younes Lakhouane \\
Site web & \url{https://wepwn.ma} \\
Siège & \acompleter[ville] \\
Année de création & \acompleter[année] \\
Effectif & \acompleter \\
\bottomrule
\end{tabularx}
\end{table}

\subsection{Cadre du stage}
Le stage s'est déroulé au sein de WePwn, sous l'encadrement professionnel de
M. Younes Lakhouane, fondateur de l'entreprise, et sous l'encadrement académique de
Mme K. Chetioui (ENSA de Fès). Il a commencé le \datedebut{}. Le sujet porte sur
\demo{}, un outil de WePwn publié en source ouverte sous licence
MIT~\cite{demo-cli-beta}, que l'entreprise destine aux équipes qui confient du travail
à des agents de code IA.

La double culture de l'entreprise explique la forme du travail demandé : construire
l'outil \emph{et} l'attaquer. Une partie du stage a consisté à développer et à
architecturer ; une autre, à chercher comment contourner ce qui venait d'être construit,
avec les méthodes du \textit{bug bounty} --- lire le code, fabriquer des commandes
piégées, fouiller les traces. Le chapitre~\ref{ch:securite} rend compte de cette seconde
partie.
